\section{Related Work}
\label{sec:related}

Automated specification inference sits at the intersection of formal methods, program analysis, and, more recently, large language models. The remainder of this section situates the present approach with respect to each of these strands.

\paragraph{From verification to inference.}
The theoretical foundations of this work lie in Hoare logic~\cite{hoare1969axiomatic} and Dijkstra's weakest-precondition calculus~\cite{dijkstra_guarded_1975, dijkstra1976discipline}. Tools such as Boogie~\cite{barnett2005weakest} and ESC/Java~\cite{flanagan2002extended, chalin2010jml} apply these formalisms to \emph{verify} user-provided specifications, and OpenJML and KeY adopt a similar verify-not-generate stance. The closest precursor on the generation side is Houdini~\cite{flanagan2001houdini}, which guesses candidate annotations from a fixed set of templates and prunes them by repeated calls to ESC/Java; the present approach extends this lineage by replacing the template guess with WP/SP-organised structural pattern matching and by extending the verification oracle to OpenJML with quantifier support. The present work inverts the verification problem, inferring contracts from code structure rather than checking given ones. Weakest-precondition calculus is a \emph{backward transformer}: it computes $\mathrm{WP}(S, Q)$ from a \emph{given} postcondition~$Q$. Specification inference has no such~$Q$---the postcondition is itself part of what must be produced---so a WP formulation presupposes half of the target artefact. The best-known gap is loop invariants (which WP requires and whose synthesis is undecidable in general), but three further considerations motivate heuristic inference. First, exact WP yields a predicate rather than a contract: its size can grow exponentially in the number of program paths, producing formulae that are neither human-readable nor the compact, idiomatic clauses (e.g.\ \verb|arr != null| or \verb|\result >= 0|) that developers write and that downstream tools consume. Second, the \emph{weakest} precondition is not the \emph{intended} specification: WP faithfully encodes implementation detail, whereas a useful contract captures designed behaviour and may be deliberately stronger and simpler---WP can return \verb|true| for a total method where the informative specification is a result--input relation. Third, sound WP requires a complete formal model of the heap, aliasing, exceptions, integer overflow, and every library call---the machinery that makes tools such as Boogie and ESC/Java heavyweight---whereas heuristic pattern matching degrades gracefully on code that no such model fully covers (unmodelled library calls, native methods), producing partial but useful specifications where exact WP produces none. This distinction matters in practice because manual specification remains the dominant cost in verification projects: in large-scale case studies it accounts for roughly 50\% of total effort~\cite{andronick_large-scale_2012, klein2014sel4, matichuk_cost_2015}. Automating the inference step targets the barrier identified by Zimmermann and Wehrheim~\cite{zimmermann2019lightweight} as the principal obstacle to lightweight specification adoption in modern development workflows.

\paragraph{Dynamic and static inference.}
Daikon~\cite{ernst2007daikon} pioneered dynamic invariant detection by generalising from execution traces; its principal limitation is its dependence on test-suite quality, because invariants are ``likely'' rather than guaranteed and coverage gaps leave specifications incomplete. PreInfer~\cite{preinfer2017} addresses this in part through symbolic analysis of path conditions collected from failing tests for C\# (built on Microsoft's Pex), but continues to require an execution infrastructure. By contrast, the approach described here operates purely statically---without test suites, execution traces, or an SMT solver---and therefore applies to any compilable Java codebase, including library code without an existing test suite. The corresponding trade-off is that heuristic pattern matching may miss specifications that would emerge from dynamic observation, particularly for complex numeric relationships.

\paragraph{Documentation-based approaches.}
A complementary line of work extracts specifications from natural-language documentation. Toradocu~\cite{toradocu2016} and its successor Jdoctor~\cite{jdoctor2018} use natural-language processing to recover exception preconditions from Javadoc comments, whereas Pandita et al.~\cite{pandita2012inferring} infer resource specifications from API documentation. Such approaches are effective when documentation is thorough but silent on undocumented methods. The complementarity is quantified in Section~\ref{sec:results}: Jdoctor captures 25 exception preconditions that JML-Inferrer misses (typically encoded in natural language), whereas JML-Inferrer captures 37 that Jdoctor misses (undocumented in Javadoc). A practical workflow benefits from combining the two.

\paragraph{LLM-based specification inference.}
Large language models have recently been applied to specification inference~\cite{ma2024specgen, endres2024can, richter2025beyondpostconditions, teuber2025nextsteps}. SpecGen~\cite{ma2024specgen} and the LLM-based approach of Endres et al.~\cite{endres2024can} prompt the model directly to emit formal contracts; Richter and Wehrheim~\cite{richter2025beyondpostconditions} extend this to full pre/post pairs and report the resulting verifier false-alarm rates. Closest in target language to the present work, Teuber and Beckert~\cite{teuber2025nextsteps} use LLMs to generate JML contracts for Java deductive verification. The comparison reported in Section~\ref{sec:results} contrasts direct LLM inference with JML-Inferrer along two axes that do not depend on a human-judged ground truth: methods covered (100\% versus 99.0\%, near-parity by construction) and run-to-run consistency (72.4\% versus 100\%, the deterministic advantage of pattern matching over sampled inference). The same comparison motivates the probe-and-backport development methodology of Section~\ref{subsec:probe-backport}: an LLM proposes candidate contracts, but a verifier filter and an agentic backport step are required before any proposal enters the engine's repertoire, because the LLM's plausibility judgement does not by itself meet the discharge bar.

\paragraph{LLM-based test generation and the oracle problem.}
LLM test generation is now an extensive subarea, recently surveyed in a systematic literature review by Zhang et al.~\cite{zhang2025llmunitsurvey} covering 115 publications. Foundational results include AthenaTest~\cite{tufano2020athenatest}, which adapted transformer architectures to unit-test generation with focal context; CodaMOSA~\cite{lemieux2023codamosa}; ChatTESTER~\cite{yuan2024chattester}; CodeT~\cite{chen2022codet}; and the empirical evaluation of Sch\"{a}fer et al.~\cite{schafer2023adaptive}. Kang, Yoon, and Yoo~\cite{kang2023libro} (LIBRO) demonstrated that LLMs can be effective few-shot testers when prompted with bug reports. The oracle problem in this setting has attracted concentrated attention. Dinella et al.~\cite{dinella2022toga} (TOGA) trained a neural model jointly to produce exceptional and assertion oracles, and Hossain and Dwyer~\cite{hossain2024togll} (TOGLL) showed that large general-purpose LLMs can produce stronger and more diverse oracles than the original neural baseline, detecting an order of magnitude more bugs. Konstantinou et al.~\cite{konstantinou2024oracles} show that LLM-generated oracles tend to encode actual rather than expected program behaviour, motivating any approach that supplies the model with an externally grounded specification. Molinelli et al.~\cite{molinelli2025oracles} evaluate 13{,}866 LLM-generated oracles on an unbiased dataset of 135 Java projects and report a mean mutation score of 43\% versus 45\% for human-written oracles, providing a calibration point for the present results. Foster et al.~\cite{foster2025meta} report industrial deployment of mutation-guided LLM test generation at Meta. Sch\"{a}fer et al.~\cite{schafer2023adaptive} report a 70\% compilation rate for LLM-generated tests; the specification-guided approach evaluated here raises this rate to 91.5\%, suggesting that specifications constrain LLM outputs toward valid test code. Rather than positioning static analysis against LLMs, the present work shows that the two are synergistic: specifications inferred by static analysis improve LLM test-generation quality beyond what signatures or source code alone provide. The TOGA/TOGLL line of work is complementary in that it learns oracles from corpora, whereas the present approach derives them from program structure; the two could in principle be ensembled. Jiang et al.~\cite{jiang2024survey} survey LLMs for code generation more broadly.

\paragraph{Mutation testing as the evaluation methodology.}
Mutation testing remains the standard means of distinguishing tests that exercise behaviour from tests that merely cover code. Papadakis et al.~\cite{papadakis2019mutation} survey the field; PIT~\cite{coles2016pitest} provides the operator set used here. Recent work has begun to study LLMs and mutation testing in combination. Wang et al.~\cite{wang2024llmmutation} present a comprehensive empirical study of LLMs for mutation testing across 851 Java bugs, and Wang et al.~\cite{wang2025mutgen} develop a mutation-guided LLM test generation technique, reporting that 100\% statement coverage can coexist with mutation scores as low as 4\%---a quantitative motivation for the mutation-based evaluation adopted in this article.

\paragraph{Specification-based testing.}
Established test-generation tools---EvoSuite~\cite{fraser2011evosuite}, Randoop~\cite{pacheco2007randoop}, and Korat~\cite{boyapati2002korat}---achieve high structural coverage but address the oracle problem only partially. Design by Contract~\cite{meyer1992design} and JML~\cite{jml, leavens2006design} address it by embedding expected behaviour in specifications, but require manual specification effort. Industrial static analysers such as Facebook Infer~\cite{calcagno2015infer} and FindBugs~\cite{hovemeyer2004finding} demonstrate the scalability of static analysis on large codebases, but focus on bug detection rather than specification generation. The structural distinction relevant to the present article is the direction of the arrow: existing JML-aware tools (including JMLUnit, EvoSuite when given JML, Randoop, and Korat) \emph{consume} formal specifications to generate tests, whereas this work \emph{infers} the specifications they consume. The two roles are complementary halves of a single pipeline; JML-Inferrer supplies the upstream piece that the established tools have always required as input. The method categorisation taxonomy adopted here extends Dragan et al.'s method stereotypes~\cite{dragan2006reverse} to identify which categories benefit most from inferred specifications.
